\section{Ba quan hệ khác nhau với chữ \enquote{tấn công}}
\label{sec:ch14-ba-quan-he}

\index{tấn công!ba quan hệ với lời giải thích|(}%
Phần mở đầu chương đã nêu rằng ba bài báo của chương đứng ở ba chỗ khác nhau so
với chữ trong nhan đề. Mục này viết ba chỗ ấy ra thành ba phát biểu, vì mỗi
phát biểu quyết định kết quả của bài tương ứng nói được về cái gì.

\index{tấn công!SHAP và LIME là nạn nhân}%
Quan hệ thứ nhất đặt lời giải thích vào vị trí nạn nhân. Một bên triển khai mô
hình có một mô hình thiên lệch mà họ muốn đưa vào dùng, và có những người kiểm
tra dùng LIME hoặc SHAP để soi mô hình ấy trước khi cho phép. Phép tấn công
dựng một mô hình mới trả về đúng dự đoán thiên lệch trên dữ liệu thật, nhưng
trả về điều gì khác khi bị LIME hoặc SHAP hỏi. Người kiểm tra nhận về một lời
giải thích sạch sẽ về một mô hình vẫn thiên lệch nguyên vẹn.
Mục~\ref{sec:ch14-phep-dung-gian} dựng lại phép ấy.

\index{tấn công!SHAP là vũ khí}%
Quan hệ thứ hai đảo chiều. Ở đó không ai lừa một lời giải thích cả; lời giải
thích là công cụ trong tay bên tấn công. Bài báo của
mục~\ref{sec:ch14-shap-lam-vu-khi} tính giá trị SHAP cho từng điểm ảnh của một
bức ảnh, đọc ra những điểm ảnh mà mô hình phân loại dựa vào nhiều nhất, rồi sửa
đúng những điểm ảnh ấy để mô hình phân loại sai. Nạn nhân là bộ phân loại, còn
SHAP đứng ở chỗ mà một phép tấn công thông thường đặt gradient. Gọi SHAP là một
mặt tấn công thì đúng với quan hệ thứ nhất và sai với quan hệ này: một mặt tấn
công là chỗ hệ thống bị chọc vào, còn ở đây SHAP là thứ dùng để chọc.

\index{tấn công!vắng mặt ở bài thứ ba}%
Quan hệ thứ ba là không có quan hệ nào. Bài báo của
mục~\ref{sec:ch14-sua-bang-nhan-qua} đề xuất một biến thể của SHAP dùng cấu
trúc nhân quả, và chỗ hỏng mà nó sửa không do ai gây ra: SHAP không phân biệt
được tương quan với nhân quả, nên khi hai đặc trưng tương quan mạnh, nó chia
điểm cho cả hai dù chỉ một trong hai thật sự tác động lên đầu ra. Chỗ hỏng ấy
có mặt trong mọi lần chạy bình thường và không cần bên tấn công nào. Chương giữ
bài ấy ở đây vì nó là phép sửa duy nhất trong danh mục nhắm vào chính cơ chế mà
quan hệ thứ nhất khai thác, chứ không vì nó nói về tấn công.

Ba quan hệ ấy không phải ba chặng của một câu chuyện, và đọc chúng như vậy là
cái bẫy mà nhan đề chương dựng sẵn. Một phép tấn công, một mặt bị tấn công, rồi
một phép sửa: bộ ba ấy nghe liền mạch tới mức người đọc dễ gán cho bài thứ ba
một bên tấn công mà nó không có, và gán cho bài thứ hai một nạn nhân mà nó không
nhắm tới.

\index{ground truth!dựng bởi bên tấn công}%
Còn một điều nữa phải nói trước khi vào bài thứ nhất, vì nó là lý do chương này
đáng đọc kỹ hơn nhan đề của nó gợi ý. Suốt từ chương~\ref{ch:do-do-trung-thuc}
tới đây, khó khăn lặp đi lặp lại là không ai có ground truth: muốn chấm điểm
một lời giải thích thì phải biết mô hình thật sự dựa vào cái gì, và với một mô
hình có sẵn thì không ai biết. Chương~\ref{ch:chi-so-khong-do-do-trung-thuc}
thoát ra bằng cách chọn những nhiệm vụ mà thiết kế của nhiệm vụ ép câu trả lời
phải là gì, và mục~\ref{sec:ch09-doi-tuong-khac-nhau} đã ghi rằng cách dựng ấy
là thứ chuyển được sang attribution đặc trưng. Phép tấn công của mục sau chuyển
nó thật, và chuyển theo một đường không ai thiết kế: bên tấn công phải tự tay
dựng mô hình thiên lệch thì mới tấn công được, nên bên ấy biết chính xác mô
hình mình dựa vào đặc trưng nào. Ground truth có sẵn, không phải vì ai đó đi
tìm nó, mà vì phép tấn công không chạy được nếu thiếu nó.

\index{tấn công!ba quan hệ với lời giải thích|)}%
